% LaTeX-ready resource-estimation section and appendix text.
% The analog-qec checkout does not currently contain a tracked manuscript .tex
% file, so this file is a drop-in source for the paper draft.

\section{Performance comparison: resource estimation}
\label{sec:resource-estimation}

The preceding simulations show that excitation-preserving simulation (EPS) can
suppress dephasing and approach the relaxation-limited coherence regime. We next
ask whether this coherence gain remains favorable once it is charged against the
space-time resources needed for a fixed observable-estimation task. We use a
50-logical-spin XY benchmark with dimensionless evolution time
\(Jt=20(2\pi)\). Taking \(t_{2\pi}=25\,\mathrm{ns}\) for a full \(2\pi\)
interaction gives
\begin{equation}
T_{\mathrm{analog}}
= \frac{Jt}{2\pi} t_{2\pi}
=0.5\,\mu\mathrm{s}.
\end{equation}
The task is to estimate final XY energy density and radial transverse
correlations from global \(X\) and global \(Y\) measurement bases. With the
bounded-estimator rule
\begin{equation}
N_{\mathrm{shots/basis}}
=\left\lceil \frac{\mathrm{Var}}{\epsilon^2}\right\rceil,
\qquad
\mathrm{Var}\le 1,\quad \epsilon=10^{-2},
\end{equation}
the comparison uses \(10^4\) shots per basis, or
\(N_{\mathrm{task}}=2\times 10^4\) circuit executions.

All points are placed on a common observable-sensitive error axis. For the
lifetime-limited branches, the coherence contribution is
\begin{equation}
H_{\mathrm{coh}}
= n_{\mathrm{eff}} T_{\mathrm{arch}}/T_{2,\mathrm{eff}},
\qquad
P_{\mathrm{err}}=1-\exp(-H),
\end{equation}
where \(n_{\mathrm{eff}}=10\) for the energy-density plus radial-correlation
task. This is a phenomenological observable-sensitivity factor: it is smaller
than the full 50-spin register-survival criterion, but larger than the local
scale relevant to energy density alone. Raw analog points use the dephasing
dominated relation
\begin{equation}
\frac{1}{T_2}=\frac{1}{2T_1}+\frac{1}{T_\phi}
\simeq \frac{1}{T_\phi},
\end{equation}
so their labels quote \(T_\phi\). EPS points quote \(T_1\) because the nominal
EPS limit is \(T_2=2T_1\). When residual pure dephasing is included, we instead
use
\begin{equation}
\frac{1}{T_{2,\mathrm{EPS,eff}}}
=\frac{1}{2T_1}
+\frac{s_\phi(\lambda,\Delta)}{T_{\phi,\mathrm{source}}}.
\label{eq:eps-residual-dephasing}
\end{equation}

The EPS lambda sweep charges the slowdown explicitly:
\begin{equation}
T_{\mathrm{EPS}}
=(1+\lambda)T_{\mathrm{analog}}.
\label{eq:eps-lambda-slowdown}
\end{equation}
The same sweep suppresses residual dephasing and coherent crosstalk through the
phenomenological envelope
\begin{equation}
s(\lambda,\Delta)
=\frac{2}{3}\left[1+(\lambda\Delta)^2\right]^{-1}
+\frac{1}{3}(1+\lambda\Delta)^{-1},
\label{eq:lambda-delta-suppression}
\end{equation}
with \(\Delta=50\) in the plotted sensitivity curves. Increasing \(\lambda\)
therefore has two competing effects: it increases task time and coherence
exposure through Eq.~\eqref{eq:eps-lambda-slowdown}, while decreasing residual
dephasing and crosstalk through Eq.~\eqref{eq:lambda-delta-suppression}.

We also include a coherent-crosstalk sensitivity term for Raw and EPS. The total
Raw/EPS exponent is
\begin{equation}
H = H_{\mathrm{coh}} + H_{\mathrm{xtalk}}.
\end{equation}
Crosstalk strengths are normalized as fractional Hamiltonian errors relative to
\(J_{\mathrm{nom}}=2\pi/t_{2\pi}\). Static crosstalk and drive-induced crosstalk
accumulate angles
\begin{equation}
\theta_{\mathrm{static}}
=\epsilon_{\mathrm{static}}J_{\mathrm{nom}}T_{\mathrm{arch}},
\qquad
\theta_{\mathrm{drive}}
=\epsilon_{\mathrm{drive}}c_{\mathrm{ctrl}}J_{\mathrm{nom}}T_{\mathrm{arch}}.
\end{equation}
For EPS, the unresolved coherent angle is suppressed at the angle level,
\begin{equation}
\theta_{\mathrm{xtalk,EPS}}
=s_{\mathrm{xtalk}}(\lambda,\Delta)
\sqrt{\theta_{\mathrm{static}}^2+\theta_{\mathrm{drive}}^2},
\end{equation}
and then mapped to a monotone twirled-envelope proxy,
\begin{equation}
p_{\mathrm{xtalk}}
=\frac{1-\exp(-\theta_{\mathrm{xtalk}}^2/2)}{2},
\qquad
H_{\mathrm{xtalk}}
=n_{\mathrm{eff,xtalk}}p_{\mathrm{xtalk}}.
\end{equation}
This term is a sensitivity model for unresolved coherent Hamiltonian error, not
a microscopic simulation of a specified crosstalk Hamiltonian.

\begin{figure}[t]
    \centering
    \includegraphics[width=\linewidth]{figures/figure_5_resource_estimation_stacked.pdf}
    \caption{Phenomenological resource-error comparison for a 50-logical-spin
    analog benchmark. The horizontal axis gives the cumulative
    observable-sensitive error exponent \(H\), with the corresponding proxy
    error probability \(P_{\mathrm{err}}=1-\exp(-H)\). The vertical axes show the
    task-level physical-space and physical-time costs after multiplying each
    circuit by the \(2\times10^4\) measurements required by the stated
    observable-estimation task. Raw points are labeled by \(T_\phi\), while EPS
    points are labeled by the relaxation-limited \(T_1\). EPS sensitivity curves
    include the lambda-dependent runtime slowdown, residual dephasing rate, and
    coherent-crosstalk envelope described in the text. STAR and surface-code
    baselines use optimistic, competitor-favorable assumptions.}
    \label{fig:resource-estimation}
\end{figure}

Figure~\ref{fig:resource-estimation} shows the resulting tradeoff. Raw analog
simulation has minimal time and space overhead, but its error exponent remains
large when the transverse lifetime is dephasing dominated and when coherent
crosstalk is not suppressed. EPS pays a small space overhead and, along the
lambda sweep, a controlled runtime slowdown. In return, it moves toward the
relaxation-limited coherence regime and suppresses the modeled residual
dephasing and crosstalk terms. The STAR and surface-code baselines are given
optimistic encoded assumptions: compact STAR blocks, no extra routing penalty,
optimistic surface-code suppression factors, and a small factory-area model.
Under these assumptions, EPS occupies a low-overhead, low-exponent regime for
this observable-estimation task.

\section{Resource-estimate model details}
\label{app:resource-estimate-details}

This appendix summarizes the calculation behind
Fig.~\ref{fig:resource-estimation}. The plotted time and space-time values are
task-level quantities. For a single circuit with architecture time
\(T_{\mathrm{arch}}\) and physical-qubit count \(n\),
\begin{equation}
T_{\mathrm{task}}=N_{\mathrm{task}}T_{\mathrm{arch}},
\qquad
nT_{\mathrm{task}}=N_{\mathrm{task}}nT_{\mathrm{arch}}.
\end{equation}
The default task uses two global measurement bases and
\(N_{\mathrm{shots/basis}}=\lceil \mathrm{Var}/\epsilon^2\rceil=10^4\), giving
\(N_{\mathrm{task}}=2\times 10^4\).

For Raw analog simulation,
\begin{equation}
T_{\mathrm{Raw}}=T_{\mathrm{analog}},
\qquad
n_{\mathrm{Raw}}=n_{\mathrm{logical}},
\qquad
H_{\mathrm{Raw}}=
n_{\mathrm{eff}}T_{\mathrm{Raw}}
\left(\frac{1}{2T_1}+\frac{1}{T_\phi}\right)
+H_{\mathrm{xtalk}}.
\end{equation}
In the plotted dephasing-dominated regime this reduces to
\(H_{\mathrm{Raw}}\simeq n_{\mathrm{eff}}T_{\mathrm{analog}}/T_\phi
+H_{\mathrm{xtalk}}\).

For EPS,
\begin{equation}
n_{\mathrm{EPS}}=2n_{\mathrm{logical}},
\qquad
T_{\mathrm{EPS}}=(1+\lambda)T_{\mathrm{analog}},
\end{equation}
and
\begin{equation}
H_{\mathrm{EPS}}
=n_{\mathrm{eff}}T_{\mathrm{EPS}}
\left(\frac{1}{2T_1}
+\frac{s_\phi(\lambda,\Delta)}{T_{\phi,\mathrm{source}}}\right)
+H_{\mathrm{xtalk}}.
\end{equation}
The crosstalk term uses the same angle-envelope model as in the main text, with
\(s_{\mathrm{xtalk}}(\lambda,\Delta)=s_\phi(\lambda,\Delta)\) for the current
sensitivity curves.

Surface-code points use an edge-colored Trotterization of the same \(Jt\), an
optimistic QEC-cycle conversion, and a logical-lifetime scaling model. STAR
points use the same Trotterized rotation count but replace Clifford+\(T\)
synthesis with injected analog-rotation resources. The STAR exponent is first
computed as a full-register rotation-plus-Clifford operation budget and then
scaled by \(n_{\mathrm{eff}}/n_{\mathrm{logical}}\) so that all schemes share
the same observable-sensitive error axis. The encoded baselines are therefore
not complete architecture-level resource counts; they are deliberately
competitor-favorable reference points for judging whether the EPS regime remains
plausibly useful under optimistic encoded assumptions.
